\scanpage{90}

\[
|\lambda_1|\ldots|\lambda_n|=\rho_1\ldots\rho_n\tag*{(24)}
\]

\deleted{En effet, le \deleted{produit}\inserted{carré du premier} mb est égal à \deleted{\unsure{$|\lambda_1$}} $|\lambda_1\ldots\lambda_n|^2=|\det u|^2=(\det u,\overline{\det u})=(\det u^*)\det u\deleted{{}^{\unsure{*}}}=\det u^*u=\rho_1^2\ldots\rho_n^2$, où $\rho_1\ldots\rho_n$ sont les valeurs propres de $(u^*u)^{1/2}$, d'où (23). \unsure{Il} \deleted{\unsure{que}} \unsure{s'ensuit alors},}

% margin note written sideways
\margin{On a \unsure{d'une part} \deleted{\unsure{d'autre} $\unsure{\lambda_1\ldots=\det}$} (25) $\boxed{\lambda_1(u)\ldots\lambda_n(u)=\det u}$ \deleted{et de} d'autre part $(\rho_1(u)\ldots\rho_n(u))^2=\rho_1(u^*u)\ldots\rho_n(u^*u)=\det(u^*u)=\det u^*\,\det u=\overline{\det u}\det u=|\det u|^2$, d'où \deleted{\unsure{(25)}} \textcircled{26} $\boxed{\rho_1(u)\ldots\rho_n(u)=|\det u|}$ En comparant (25) et (26), \unsure{tr.} $\Rightarrow$ (24)}

quand $u$ est inversible, \deleted{que} donc \deleted{\unsure{le} \unsure{tout} \unsure{par}} les $\lambda_i$ et les $\rho_i$ non nuls \deleted{que ce sont}, que les conditions du \deleted{corollaire 1 du lemme fondamental 2 sont vérifiées, dans l'hypothèse b)} \inserted{th. 1, 2°, hypo. b sont vérifiées}. Si on suppose que $f$ est de plus \uline{définie} et continue pour $r_i\geq0$, comme on aura (23) pour \deleted{tout} $u$ inversible et que les deux mbs sont des fonctions continues de $u$ comme on vérifie aussitôt, il s'ensuit que l'inégalité sera vérifiée pour \deleted{tout} $u$. \unsure{Il|Le}

\deleted{\uline{Corollaire 5}} résulte p. ex du corollaire 4 que si $u$ est inversible, les formules (21) sont vérifiées même si $p$ n'est pas positif. Notons que \deleted{sous les conditions du corollaire 3, on tire de (24) et des inégalités $|\lambda_1\ldots\lambda_k|\leq\rho_1\ldots\rho_k$, les relations analogues}

% slanted
\margin{\deleted{\unsure{(25)|(24)} $|\lambda_k|\unsure{\leq|\geq}\rho_{\unsure{n}}\ldots\rho_n$ pour $(\unsure{1\leq k\leq n})$, ce qui \illegible{2} pour $\unsure{\rho_k}\ldots\unsure{\rho_n}$, \illegible{3} \unsure{la} \illegible{2}} \textcircled{?}}

\uline{Corollaire 5} Soit $u$ un opérateur compact dans un espace de Hilbert $E$, soit $N$ une norme de Schatten sur $R^{\deleted{\unsure{I}}k}$, \deleted{où $I=\dim(1,\ldots n)$ soit $\dim E=n$, et $I=$ ensemble des entiers si} \inserted{où $k\leq\dim E$. Alors} on a

\[
N(|\lambda_1|,\ldots,|\lambda_k|)\leq N(\rho_1,\ldots\rho_k)\tag*{(27)}
\]

\deleted{D'où, \unsure{$N($} \illegible{1}} Par suite, si la dimension de $E$ est infinie et si $N$ est une norme de Schatten sur \deleted{\unsure{tout}} l'espace $R^{(\underline N)}$, \deleted{\unsure{lorsque} $\underline N$ \unsure{désigne}} \inserted{on a}

\[
N((|\lambda_i|))\leq N((\rho_i))\tag*{\unsure{(28)|(26)}}
\]

Dans cet énoncé, $\underline N$ désigne l'ensemble des entiers, donc $R^{(\underline N)}$ l'espace des suites de \unsure{nbs entrs|nbs réels|nbs scalres} \deleted{qui sont} dont seul un nb fini de termes sont non nuls. Une norme sur $R^{(\underline N)}$ est dite une norme de Schatten, si elle induit une norme de Schatten sur \deleted{tout} un espace $R^n$ (plongé de façon naturelle dans $R^{(\underline N)}$).

% archive sheet number: 91
